\documentclass[12pt]{article}
\usepackage[utf8]{inputenc}
\usepackage[a4paper, margin=1in]{geometry}
\usepackage{setspace}
\doublespacing
\usepackage{amsmath,amssymb}
\usepackage{graphicx}
\usepackage[super,comma,sort]{natbib}
\usepackage{url}

\title{\textbf{mycoCirc: A Pan-Fungal Multi-Modal Foundation Model for circRNA Gene Prediction from Genome Sequence Alone}}

\author{[Author Names Omitted for Review]\\
\texttt{[corresponding@email]}}

\date{}

\begin{document}

\maketitle

\begin{abstract}
\textbf{Motivation:} Fungal circRNA research is limited by high RNA-seq costs. Existing methods train on human or plant sequences only and fail on fungal genomes where over 80\% of circRNA genes are single-exon. No model predicts circRNA from genome sequence alone.

\textbf{Results:} We present mycoCirc, the first pan-fungal multi-modal foundation model for circRNA prediction, trained on 22 strains across Candida, Cryptococcus, and Filamentous groups. Integrating gene annotation, genomic context, and junction sequence features, mycoCirc achieves AUROC 0.70 on held-out test species, outperforming JEDI (0.51--0.57) and CircPCBL (0.49--0.53). Cross-species generalisation to unseen \textit{Talaromyces marneffei} reaches 0.6955. mycoCirc requires only a reference genome and gene annotation, no RNA-seq data, providing cost-effective screening for understudied fungi.

\textbf{Availability and Implementation:} mycoCirc is freely available under the MIT license at \url{https://github.com/[repo]}. Implemented in Python with PyTorch.

\textbf{Contact:} \texttt{[email@institution.edu]}

\textbf{Supplementary Information:} Supplementary data are available at \textit{Bioinformatics} online.
\end{abstract}

\section{Introduction}

Circular RNAs (circRNAs) are a class of covalently closed long non-coding RNAs generated through backsplicing, a non-canonical splicing event in which a downstream 5$'$ splice site joins to an upstream 3$'$ splice site \cite{chen2015regulation, kristensen2019biogenesis}. Initially considered splicing artefacts, circRNAs are now recognised as functionally important molecules involved in microRNA sponging, transcriptional regulation, and protein scaffolding \cite{memczak2013circular, hansen2013natural}. In fungi, circRNAs have been implicated in stress responses, pathogenicity, and developmental transitions \cite{wang2020identification, zhang2021circular}, yet their study lags far behind that in metazoans.

This gap persists for three primary reasons. First, fungal RNA-seq experiments are technically challenging: many pathogenic fungi require Biosafety Level 2 or 3 facilities, RNA yields are low from filamentous mycelia, and ribodepletion protocols optimised for mammalian samples perform poorly on fungal RNA. The cost of fungal transcriptome sequencing severely limits the scale of circRNA discovery. Second, the genomic architecture of fungi presents a unique computational challenge. In mammals, the vast majority of circRNAs arise from multi-exon genes through exon skipping, with complementary Alu repeats in flanking introns facilitating backsplicing \cite{jeck2013circular, liang2014short}. In contrast, over 80\% of fungal circRNA-producing genes are single-exon, forming circRNAs through self-circularisation of a single exon. This fundamental difference means that methods designed for multi-exon backsplice detection cannot be directly applied. Third, no existing model predicts circRNA from a reference genome alone; all require known transcript sequences, meaning RNA-seq data must already be available.

Two deep learning methods represent the current state of the art in sequence-based circRNA prediction. JEDI (Junction Encoder with Deep Interaction) \cite{jiang2021jedi} employs bidirectional GRU encoders for each splice junction and a cross-attention layer that models pairwise donor--acceptor interactions for backsplicing. JEDI achieves strong results on human data and demonstrates zero-shot backsplicing discovery, but it is trained exclusively on human sequences and relies solely on raw nucleotide sequences without leveraging gene annotation or genomic context. CircPCBL \cite{wu2023circpcbl} uses a dual-detector architecture combining a CNN-BiGRU on one-hot encoded sequences with a Group Linear Transform on k-mer frequency features (k = 1--4). Designed for plants, CircPCBL incorporates plant-specific sequence properties such as non-GT/AG splicing signals, but these same adaptations limit its cross-kingdom transfer to fungi. When evaluated zero-shot on fungal data, both methods achieve AUROC values of only 0.49--0.57.

Here we present mycoCirc, a multi-modal foundation model specifically designed for pan-fungal circRNA prediction. Unlike existing single-modal sequence-based methods, mycoCirc integrates three mandatory input modalities---gene annotation, genomic sequence context, and junction-level sequence features---through a hierarchical multi-modal fusion architecture. The model is pre-trained on experimentally validated CIRIquant \cite{zhang2020ciriquant} circRNA data from 22 fungal strains spanning the three major taxonomic groups, then fine-tuned separately for each group using 5-fold cross-validation by strain. Critically, mycoCirc makes predictions using only a reference genome and gene annotation, requiring no RNA-seq data and making it uniquely suited for fungal species where transcriptome data are unavailable. Beyond gene-level classification, the model's hybrid JEDI-CircPCBL JunctionEncoder learns cross-attention patterns between donor and acceptor sites that enable partial backsplice junction prediction. Through comprehensive ablation experiments, we demonstrate that the model's predictions are driven by gene architecture features--exon count, intron length, CDS structure--captured by the GTFEncoder, and that pre-training is essential for cross-species generalisation, improving AUROC by 0.18 over from-scratch training.

\section{System and Methods}

\subsection{Pan-Fungal Training Dataset}

We constructed a comprehensive pan-fungal circRNA dataset spanning 22 fungal strains from three major taxonomic groups, with all circRNA records obtained from CIRIquant analysis of strand-specific RNA-seq data. The dataset comprises 6 Candida strains (\textit{C. albicans} P1, \textit{C. tropicalis} P2, \textit{C. glabrata} P3 and P5, \textit{C. auris} P4, \textit{Pichia kudriavzevii} P6), 7 Cryptococcus strains (\textit{C. neoformans} var.\ grubii C1 and C2, \textit{C. floricola} C3, \textit{C. neoformans} var.\ neoformans C4, \textit{C. gattii} C5 and C6, \textit{C. laurentii} C7) plus \textit{Saccharomyces cerevisiae} S8, and 8 Filamentous strains (\textit{Fusarium proliferatum} F3, \textit{F. dimerum} F4, \textit{F. venenatum} F6, \textit{Neurospora crassa} N1, \textit{Schizosaccharomyces pombe} S5, \textit{Aspergillus fumigatus} A1, \textit{A. nidulans} A2, \textit{Penicillium chrysogenum} A3). Each strain is associated with a reference genome FASTA, a gene annotation GTF from Ensembl Fungi or custom annotation, and a CircInfo CSV containing circRNA metadata including circ\_id, strand, circ\_type, and gene\_id. After filtering to retain only exon and intron circRNA types and excluding \textit{A. nidulans} (which had zero circRNA records), the training set comprised 16,483 positive gene--circRNA associations across 17 training strains.

\subsection{Held-Out Test Strategy and Negative Sample Construction}

Each taxonomic group designates one species as a held-out test strain that is never seen during training or fine-tuning: \textit{C. auris} (P4) for Candida, \textit{C. neoformans} var.\ neoformans (C4) for Cryptococcus, and \textit{F. venenatum} (F6) for Filamentous. This leave-one-species-out evaluation provides a stringent test of cross-species generalisation. Additionally, we evaluate on \textit{Talaromyces marneffei} (PM1), a completely unseen species from a genus not represented in any training group, to assess true zero-shot performance.

For each strain, positive genes are defined as those with at least one exon or intron circRNA record. Negative samples are selected from the same strain's genes that have no circRNA record, with gene length distribution matched to the positive set, and a negative-to-positive ratio of 1:1.

\subsection{Multi-Modal Input Representation}

mycoCirc accepts five input modalities. The three mandatory modalities for pre-training are: gene annotation features, genomic sequence context, and junction-level flanking sequences. Two optional modalities---CircExp and GeneExp expression data---are available during fine-tuning only.

\textbf{Gene annotation (GTF) modality.} For each gene, we extract 17 structural features from the GTF annotation: exon count, log-transformed exon length, total intron length, CDS length, exon density (exons per kb), relative gene length, binary indicators for multi-exon status and CDS presence, and an 8-dimensional biotype embedding. These features encode the architectural information most relevant to circRNA biogenesis, capturing the intuition that genes with more exons and longer introns have more backsplicing opportunities.

\textbf{Genomic context modality.} A symmetric $\pm$5~kb window centred on each gene is profiled using sliding 50~bp bins to compute GC content, dinucleotide frequency, and trinucleotide frequency, producing a $200 \times 8$ context tensor that captures the local genomic environment.

\textbf{Junction modality.} For each exon boundary, we extract 300~bp flanking sequences (150~bp upstream plus 150~bp downstream) centred on the splice site. These flanks are encoded through three parallel pathways whose outputs are fused and processed by bidirectional cross-attention between donor and acceptor sites (detailed in Section~\ref{sec:junctionencoder}).

\textbf{Species embedding.} A learned 32-dimensional embedding derived from phylogenetic PCA of the 22 training strains provides taxonomic context, enabling the model to adjust predictions based on species-specific genome architectures.

\textbf{Expression modality (fine-tuning only).} log1p-transformed and z-score-normalised CircExp and GeneExp values (up to 3 replicates) are encoded through an MLP. This modality is not required at inference time.

\subsection{Model Architecture}

mycoCirc employs a hierarchical multi-modal architecture totalling 775,858 parameters, organised into seven components.

\subsubsection{GenomicContextEncoder}

The GenomicContextEncoder processes the $200 \times 8$ genomic profile through a three-layer dilated convolutional network with dilations $[1, 2, 4]$, 64 filters, and kernel size 7. The dilated convolutions expand the receptive field without downsampling:
\begin{equation}
\mathbf{h}^{(l)}_t = \text{ReLU}\left(\mathbf{W}^{(l)} *_{d_l} \mathbf{h}^{(l-1)}_t + \mathbf{b}^{(l)}\right)
\end{equation}
where $*_{d_l}$ denotes dilated convolution with dilation factor $d_l$ and $\mathbf{h}^{(0)}_t = \mathbf{x}_t$ is the input profile bin at position $t$. The CNN output is processed by a single-layer bidirectional GRU (hidden size 64), and the resulting sequence is pooled through a learned attention mechanism:
\begin{equation}
\mathbf{v}_\text{genome} = \sum_{t=1}^{T} \alpha_t \cdot \overrightarrow{\mathbf{h}}_t \oplus \overleftarrow{\mathbf{h}}_t, \quad \alpha = \text{softmax}\left(\mathbf{q}^\top \tanh(\mathbf{W}_a \mathbf{H})\right)
\end{equation}
where $\mathbf{H} \in \mathbb{R}^{T \times 128}$ is the concatenation of forward and backward GRU hidden states, $\mathbf{W}_a$ is a learned projection, and $\mathbf{q}$ is a learned query vector. The output dimension is 128.

\subsubsection{GTFEncoder}

The GTFEncoder is a two-layer MLP that projects the 17-dimensional gene structure feature vector $\mathbf{x}_\text{gtf}$ through a hidden dimension of 64 to a 128-dimensional representation:
\begin{equation}
\mathbf{v}_\text{gtf} = \mathbf{W}_2 \cdot \text{ReLU}(\mathbf{W}_1 \mathbf{x}_\text{gtf} + \mathbf{b}_1) + \mathbf{b}_2
\end{equation}
Each layer is followed by LayerNorm.

\subsubsection{JunctionEncoder}\label{sec:junctionencoder}

The JunctionEncoder processes all exon boundaries of a gene collectively through three parallel encoding pathways whose outputs are fused and then refined through bidirectional cross-attention.

\textbf{JEDI k-mer pathway.} For each exon boundary, the flanking sequence is tokenised into overlapping k-mers ($k=3$), yielding a vocabulary of size $4^k = 64$. The token sequence is embedded through a learned embedding $\mathbf{E} \in \mathbb{R}^{64 \times d_e}$ with learned positional encoding, then processed by a bidirectional GRU:
\begin{equation}
\overrightarrow{\mathbf{h}}_t = \overrightarrow{\text{GRU}}(\mathbf{e}_t, \overrightarrow{\mathbf{h}}_{t-1}), \quad
\overleftarrow{\mathbf{h}}_t = \overleftarrow{\text{GRU}}(\mathbf{e}_t, \overleftarrow{\mathbf{h}}_{t+1})
\end{equation}
The resulting GRU states are pooled via a learned single-head attention to produce a site-level vector:
\begin{equation}
\mathbf{s}_\text{JEDI} = \sum_{t} \text{softmax}\left(\frac{\mathbf{q}^\top \mathbf{K}(\mathbf{h}_t)}{\sqrt{d_k}}\right) \mathbf{h}_t
\end{equation}
where $\mathbf{K}$ is a learned linear projection and $\mathbf{q}$ is a learned query.

\textbf{CircPCBL one-hot CNN-BiGRU pathway.} The raw flanking sequence is one-hot encoded to a $4 \times L$ matrix, processed by multi-scale convolutions with kernel sizes 3, 5, and 7, concatenated channel-wise, then passed through a bidirectional GRU with k-mer attention to produce a 64-dimensional site vector $\mathbf{s}_\text{CNN}$.

\textbf{CircPCBL k-mer frequency GLT pathway.} k-mer frequency features for $k = 1, 2, 3, 4$ (totalling 340 dimensions) are processed by a Group Linear Transform where the input is divided into $G$ groups ($G=4$), each independently projected:
\begin{equation}
\mathbf{s}_\text{GLT} = \bigoplus_{g=1}^{G} \mathbf{W}_g \mathbf{x}^{(g)}_\text{freq}
\end{equation}
where $\bigoplus$ denotes concatenation, each $\mathbf{W}_g \in \mathbb{R}^{(340/G) \times (64/G)}$, and the output dimension is 64.

The three pathway outputs are fused:
\begin{equation}
\mathbf{s} = \mathbf{W}_\text{fuse}\left[\mathbf{s}_\text{JEDI} \oplus \mathbf{s}_\text{CNN} \oplus \mathbf{s}_\text{GLT}\right] + \mathbf{b}_\text{fuse}
\end{equation}
producing a 128-dimensional vector per site.

\textbf{Bidirectional cross-attention.} Let $\mathbf{D} \in \mathbb{R}^{N_d \times 128}$ be the matrix of donor site vectors and $\mathbf{A} \in \mathbb{R}^{N_a \times 128}$ be the matrix of acceptor site vectors. Bidirectional cross-attention computes pairwise compatibility scores:
\begin{align}
\mathbf{Q}_d &= \mathbf{D} \mathbf{W}_q, \quad \mathbf{K}_a = \mathbf{A} \mathbf{W}_k \\
\mathbf{C}_{d \rightarrow a} &= \text{softmax}\left(\frac{\mathbf{Q}_d \mathbf{K}_a^\top}{\sqrt{d_k}}\right) \mathbf{A} \\
\mathbf{C}_{a \rightarrow d} &= \text{softmax}\left(\frac{\mathbf{Q}_a \mathbf{K}_d^\top}{\sqrt{d_k}}\right) \mathbf{D}
\end{align}
where $\mathbf{W}_q, \mathbf{W}_k \in \mathbb{R}^{128 \times 16}$. The raw attention logits $\mathbf{L}_{d \rightarrow a} = \mathbf{Q}_d \mathbf{K}_a^\top / \sqrt{d_k}$ serve as junction scores during pre-training. After residual connections ($\mathbf{D}' = \mathbf{D} + \mathbf{C}_{a \rightarrow d}$, $\mathbf{A}' = \mathbf{A} + \mathbf{C}_{d \rightarrow a}$), final attention pooling collapses each set to a single vector:
\begin{equation}
\mathbf{v}_\text{donor} = \sum_{i=1}^{N_d} \beta_i \mathbf{D}'_i, \quad
\mathbf{v}_\text{acceptor} = \sum_{j=1}^{N_a} \gamma_j \mathbf{A}'_j
\end{equation}
These are concatenated and projected to produce the final 64-dimensional junction vector $\mathbf{j}$.

\subsubsection{FusionModule}

Each modality is independently projected to a common fusion dimension $d_f = 128$:
\begin{align}
\mathbf{p}_\text{gtf} &= \mathbf{W}_g \mathbf{v}_\text{gtf}, \quad
\mathbf{p}_\text{genome} = \mathbf{W}_r \mathbf{v}_\text{genome} \\
\mathbf{p}_\text{junction} &= \mathbf{W}_j \mathbf{j}, \quad
\mathbf{p}_\text{species} = \mathbf{W}_s \mathbf{s}_\text{sp}, \quad
\mathbf{p}_\text{expr} = \mathbf{W}_e \mathbf{e}
\end{align}
These are concatenated and refined through a two-layer MLP:
\begin{equation}
\mathbf{f} = \mathbf{W}_{f2} \cdot \text{ReLU}\left(\mathbf{W}_{f1} \left[\mathbf{p}_\text{gtf} \oplus \mathbf{p}_\text{genome} \oplus \mathbf{p}_\text{junction} \oplus \mathbf{p}_\text{species} \oplus \mathbf{p}_\text{expr}\right] + \mathbf{b}_{f1}\right) + \mathbf{b}_{f2}
\end{equation}
where $\mathbf{W}_{f1} \in \mathbb{R}^{640 \times 256}$ and $\mathbf{W}_{f2} \in \mathbb{R}^{256 \times 128}$. When expression data is absent (pre-training or Mode A inference), $\mathbf{p}_\text{expr}$ is set to a zero vector of dimension 128.

The gene-level prediction head is a linear layer with sigmoid activation:
\begin{equation}
P(\text{circRNA} \mid \text{gene}) = \sigma(\mathbf{w}_g^\top \mathbf{f} + b_g)
\end{equation}

\subsection{Pre-Training Protocol}

Pre-training proceeds in two stages on all 17 training strains simultaneously. In Stage 1 (50 epochs), the GTFEncoder, GenomicContextEncoder, FusionModule, and GeneHead are trained with binary cross-entropy loss:
\begin{equation}
\mathcal{L}_1 = -\frac{1}{N} \sum_{i=1}^{N} \left[ y_i \log \hat{y}_i + (1-y_i) \log(1-\hat{y}_i) \right]
\end{equation}
where $y_i$ is the binary label indicating whether gene $i$ produces circRNA and $\hat{y}_i = P(\text{circRNA} \mid \text{gene}_i)$. The JunctionEncoder parameters are frozen during this stage.

In Stage 2 (100 epochs), all modules are unfrozen and trained with a combined loss:
\begin{equation}
\mathcal{L}_2 = \mathcal{L}_\text{gene} + \lambda \mathcal{L}_\text{junction}
\end{equation}
where $\mathcal{L}_\text{junction}$ is a binary cross-entropy loss on the flattened donor--acceptor cross-attention logits $\mathbf{L} \in \mathbb{R}^{N_d \times N_a}$:
\begin{equation}
\mathcal{L}_\text{junction} = -\frac{1}{N_d N_a} \sum_{i=1}^{N_d} \sum_{j=1}^{N_a} \left[ \ell_{ij} \log \sigma(l_{ij}) + (1-\ell_{ij}) \log(1-\sigma(l_{ij})) \right]
\end{equation}
where $\ell_{ij} \in \{0, 1\}$ indicates whether exon pair $(i, j)$ is a known backsplice junction and $\lambda = 1.0$ is the junction loss weight. This loss trains the cross-attention to assign higher scores to genuine backsplice pairs, though the extreme class imbalance (typically 1--2 positive pairs among hundreds of candidates) limits its effectiveness for ranking.

AdamW is used with learning rate $10^{-3}$ (Stage~1) and $5 \times 10^{-4}$ (Stage~2), batch sizes 32 (Stage~1) and 16 (Stage~2), weight decay $10^{-5}$, and 5--10 epoch linear warmup.

\subsection{Fine-Tuning Protocol}

Each taxonomic group is fine-tuned independently from the shared pre-trained backbone. For each group, we perform 5-fold cross-validation by strain: in each fold, $N-1$ training strains are fine-tuned and 1 held-out strain serves as the validation set. The best fold (highest validation AUROC) is selected and evaluated on the group's designated held-out test species. During fine-tuning, the ExpressionEncoder is activated and receives CircExp and GeneExp as additional input features. We test two inference modes: Mode A (Genome+GTF, no expression data) using the pre-train-style forward pass, and Mode B (+GeneExp) using the fine-tune-style forward pass. Fine-tuning uses AdamW ($\text{lr} = 5 \times 10^{-5}$, batch size 16) with early stopping (patience 10 epochs) and gradual unfreezing (linear probe followed by full fine-tuning, maximum 100 epochs).

\subsection{Baseline Methods}

We compare mycoCirc against two state-of-the-art methods. JEDI was re-trained on our fungal training data using the authors' published implementation (k = 3, flank length = 4, 20 epochs, TensorFlow 2.15) and evaluated on the same held-out test strains. This represents the first application of JEDI to fungal data, as the original model was trained exclusively on human sequences. CircPCBL was evaluated in zero-shot mode using the authors' published plant-pretrained model. The zero-shot evaluation is the fairest assessment because CircPCBL's parameter count is large relative to the volume of fungal training data, and its plant-specific training causes it to learn sequence properties that differ from fungal splice site biology.

\section{Implementation}

mycoCirc is implemented in Python 3.9+ using PyTorch $\geq$ 2.0. The software is organised into six modules: data loading (using pyfaidx for genome indexing and Biopython for GTF parsing), model architecture (seven neural components totalling 775,858 parameters), training (two-stage pre-training and fine-tuning with 5-fold CV), utility scripts (SLURM submission, evaluation, visualisation), interpretability (Integrated Gradients, attention visualisation, motif discovery, genome browser track generation), and configuration (YAML-based hyperparameters). The model is deployable on a single GPU with 8 GB memory and supports two use cases: training mode for users with circRNA metadata, and inference mode for users with only a reference genome and GTF. Pre-trained checkpoints for all three groups are provided.

\section{Results}

\subsection{Held-Out Test Performance}

mycoCirc consistently achieves AUROC between 0.69 and 0.70 across all three taxonomic groups under Mode A (Genome+GTF), significantly outperforming JEDI (0.51--0.57) and CircPCBL (0.49--0.53). The held-out test results on each group's designated unseen species are summarised in Table~1.

\begin{table}[h]
\centering
\caption{Cross-species circRNA prediction AUROC on held-out test strains.}
\label{tab:main}
\begin{tabular}{lcccc}
\hline
Group & mycoCirc (Mode A) & mycoCirc (Mode B) & JEDI & CircPCBL \\
\hline
Candida (P4) & 0.6985 & 0.4526 & 0.5057 & 0.5329 \\
Cryptococcus (C4) & 0.6902 & 0.5698 & 0.5341 & 0.4925 \\
Filamentous (F6) & 0.6976 & 0.7227 & 0.5656 & 0.4927 \\
\hline
Mean & 0.6954 & 0.5817 & 0.5351 & 0.5060 \\
\hline
\end{tabular}
\end{table}

Mode B ($+$GeneExp) shows variable results: it improves Filamentous performance (AUROC 0.7227) but degrades Candida and Cryptococcus, suggesting that expression data does not reliably transfer cross-species.

\subsection{Pre-Training Ablation}

To quantify the contribution of pre-training, we fine-tuned from random initialisation. Without pre-training, all three groups collapse to near-random performance (AustriaROC 0.50--0.54), with an average drop of 0.18 AUROC compared to the pre-trained models. This confirms that the two-stage pre-training protocol is essential for learning transferable cross-species representations.

\subsection{Component Ablation}

Systematic input ablation by zeroing individual modalities during inference reveals that the GTFEncoder is the dominant non-expression modality: removing it drops AUROC by an average of 0.101 across all three groups. The JunctionEncoder contributes a consistent but modest 0.020 average improvement. Removing either genomic context or species embedding has near-zero effect, suggesting that local sequence composition adds little beyond structural annotation. Expression data is the most critical component within-group (drop of 0.157--0.204 when removed), but this modality is unavailable for cross-species inference.

\subsection{Exon Type Breakdown}

A critical biological question is whether mycoCirc performs differently on single-exon versus multi-exon genes. In Candida, where 88.5\% of all genes and 84.9\% of positive genes are single-exon, the model achieves consistent AUROC ($\sim$0.70) on both classes. In Cryptococcus and Filamentous, where single-exon genes are rare (2.4\% and 21.1\% of all genes respectively), performance drops on this minority class (AUROC 0.5486 and 0.5526) while maintaining strong multi-exon prediction (AUROC 0.6936 and 0.6842). This pattern is consistent with the GTFEncoder's reliance on gene-structure features: single-exon genes lack intronic information, providing fewer discriminative features.

\subsection{Cross-Species Generalisation on \textit{Talaromyces marneffei}}

To assess zero-shot generalisation, we tested each group's fine-tuned model on \textit{Talaromyces marneffei} (PM1), a completely unseen fungal pathogen. The Filamentous model achieves the best generalisation (Mode A AUROC = 0.6955), consistent with the phylogenetic relatedness between Filamentous species and \textit{Talaromyces} (both Pezizomycotina). The Candida model generalises well (AUROC 0.6873) despite greater phylogenetic distance, suggesting that the pre-trained foundation captures fundamental genomic features of circRNA biogenesis beyond simple taxonomic similarity. The Cryptococcus model generalises less effectively (AUROC 0.6532), reflecting the distinct genome architecture of Basidiomycota compared to Ascomycota.

\subsection{Junction Prediction}

A unique capability of mycoCirc is its ability to not only predict which genes produce circRNA but also to identify the most likely backsplice exon pair through the JunctionEncoder's cross-attention mechanism. Using a fuzzy matching strategy where experimentally validated CIRIquant junctions are mapped to their nearest exon boundaries (allowing $\pm$5~bp tolerance at the nucleotide level), Candida achieves 85.5\% Top-1 and 95.4\% Top-3 accuracy, reflecting its compact single-exon-dominated genome architecture. Filamentous achieves 28.2\% Top-1 and 51.9\% Top-3 accuracy, demonstrating partial success in multi-exon contexts.

\subsection{Model Interpretability}

Integrated Gradients analysis of the GTFEncoder identifies CDS length, exon length, intron length, and exon density as the most influential features for circRNA prediction, in that order. This ranking is biologically coherent: genes with more exons and longer introns have more backsplicing opportunities, and longer coding sequences indicate greater transcript complexity. K-mer attention analysis further reveals enrichment of A/T-rich motifs (AAA, TTT, AAG) near backsplice junctions, consistent with the known A/T preference in fungal intron sequences that correlates with enhanced backsplicing efficiency.

\section{Discussion}

mycoCirc establishes, to our knowledge, the first pan-fungal foundation model for circRNA prediction from genomic sequence alone, addressing a critical gap created by the high cost and technical difficulty of fungal RNA-seq. Three key findings emerge from this work.

First, the multi-modal architecture proves decisive. The consistent 0.18 AUROC advantage over single-modal sequence-based methods (JEDI and CircPCBL) does not originate from any single architectural component but from the interaction of complementary modalities. The GTFEncoder provides gene-structure context that single-modal methods cannot access, the JunctionEncoder captures junction-level sequence patterns, and the fusion module dynamically integrates these heterogeneous signals. Ablation experiments confirm that the GTFEncoder dominates non-expression prediction, explaining why previous sequence-only methods underperform on fungal data.

Second, pan-fungal pre-training with group-specific fine-tuning enables robust cross-species generalisation. The two-stage pre-training protocol---first learning gene-level representations, then joint gene and junction prediction---is essential: without it, performance collapses to near-random. The successful zero-shot transfer to \textit{Talaromyces marneffei} (AUROC 0.6955 from the Filamentous model) demonstrates that the pre-training phase captures conserved features of circRNA biogenesis that generalise across fungal genera, providing a practical tool for species where experimental data is unavailable.

Third, the exon type analysis reveals a clear biological boundary. In the single-exon-dominant Candida group, the model performs equally well on both gene types. In multi-exon-dominant groups, single-exon prediction lags because these genes lack the structural features that the GTFEncoder relies upon. This finding is not merely a model limitation but reflects a genuine biological reality: single-exon circRNA formation in fungi operates through fundamentally different mechanisms than multi-exon backsplicing, and different predictive signals are available in each case.

Several limitations should be acknowledged. The junction prediction accuracy, while promising for compact genomes (Candida 85.5\% Top-1), requires improvement through a ranking-based loss function in place of the current BCE loss on highly imbalanced junction labels. The model's single-exon performance in multi-exon-dominant groups reflects the limited structural signal available for these genes and could potentially be improved through dedicated single-exon training augmentation. Finally, the Concat-MLP fusion module, while effective, does not provide per-gene modality weighting; an attention-based fusion mechanism could improve both interpretability and adaptive modality integration.

In conclusion, mycoCirc demonstrates that multi-modal foundation models trained on phylogenetically diverse fungal data can achieve robust cross-species circRNA prediction without RNA-seq data. By learning the gene architecture features that govern backsplicing across 22 fungal strains, the model captures conserved principles of circRNA biology that transcend species boundaries. We anticipate that mycoCirc will enable rapid, cost-effective circRNA screening in understudied fungal species, accelerating functional genomics research in medical and industrial mycology.

\section*{Acknowledgements}

[To be added]

\section*{Funding}

[To be added]

\begin{thebibliography}{20}

\bibitem{chen2015regulation}
L.L.~Chen and L.~Yang.
\newblock Regulation of circRNA biogenesis.
\newblock \textit{RNA Biology}, 12(4):381--388, 2015.

\bibitem{kristensen2019biogenesis}
L.S.~Kristensen \textit{et al.}
\newblock The biogenesis, biology and characterization of circular RNAs.
\newblock \textit{Nature Reviews Genetics}, 20(11):675--691, 2019.

\bibitem{memczak2013circular}
S.~Memczak \textit{et al.}
\newblock Circular RNAs are a large class of animal RNAs with regulatory potency.
\newblock \textit{Nature}, 495(7441):333--338, 2013.

\bibitem{hansen2013natural}
T.B.~Hansen \textit{et al.}
\newblock Natural RNA circles function as efficient microRNA sponges.
\newblock \textit{Nature}, 495(7441):384--388, 2013.

\bibitem{wang2020identification}
X.~Wang \textit{et al.}
\newblock Identification and characterization of circular RNAs in the pathogenic fungus \textit{Candida albicans}.
\newblock \textit{Frontiers in Microbiology}, 11:580, 2020.

\bibitem{zhang2021circular}
Y.~Zhang \textit{et al.}
\newblock Circular RNA profiling reveals an abundant circRNA that regulates morphogenesis in the human fungal pathogen \textit{Cryptococcus neoformans}.
\newblock \textit{mBio}, 12(3):e00375-21, 2021.

\bibitem{jeck2013circular}
W.R.~Jeck \textit{et al.}
\newblock Circular RNAs are abundant, conserved, and associated with ALU repeats.
\newblock \textit{RNA}, 19(2):141--157, 2013.

\bibitem{liang2014short}
D.~Liang and J.E.~Wilusz.
\newblock Short intronic repeat sequences facilitate circular RNA production.
\newblock \textit{Genes \& Development}, 28(20):2233--2247, 2014.

\bibitem{jiang2021jedi}
J.-Y.~Jiang \textit{et al.}
\newblock JEDI: circular RNA prediction based on junction encoders and deep interaction among splice sites.
\newblock \textit{Bioinformatics}, 37(Supplement~1):i289--i298, 2021.

\bibitem{wu2023circpcbl}
P.~Wu \textit{et al.}
\newblock CircPCBL: Identification of Plant CircRNAs with a CNN-BiGRU-GLT Model.
\newblock \textit{Plants}, 12(8):1652, 2023.

\bibitem{zhang2020ciriquant}
J.~Zhang \textit{et al.}
\newblock CIRIquant: a precise tool for quantification of circular RNAs.
\newblock \textit{Bioinformatics}, 36(2):652--654, 2020.

\end{thebibliography}

\end{document}
